\section{解析延拓与 Fourier 衰减}
\label{sec:08-Numerical-Analysis/09-Complex-Analysis-and-Residues/01}

在 \([-\pi,\pi]\) 上取两个周期函数做 Fourier 展开。一个是 \(f(x)=1/(2-\cos x)\)，另一个是把 Runge 函数搬成周期形式的 \(g(x)=1/(1+25\sin^2 x)\)。两条曲线都无穷次可微，都有界，画出来光滑得挑不出毛病。

截到 20 项，\(f\) 的截断误差已经落在 \(10^{-11}\) 附近；\(g\) 还停在 \(10^{-2}\)。想让 \(g\) 也降到 \(10^{-11}\)，得取到一百三十项上下。

同样的算法，同样的截断方式，差出六七倍的项数。而实轴上找不到任何解释：光滑性一样，有界性一样，周期一样。\hyperref[sec:08-Numerical-Analysis/08-Fourier-and-Spectral-Methods/01]{``从 Fourier 级数到连续变换''}给过一条粗线索——光滑函数的系数按幂律衰减，解析函数按几何衰减。可 \(f\) 与 \(g\) 都解析，都在几何衰减，差的是那个衰减率。\(e^{-a\abs{k}}\) 里的 \(a\) 究竟由什么定？实轴上没有写。

\subsection{把两条曲线抬进复平面，差别立刻显形}

\(f\) 的分母在 \(\cos z=2\) 处为零。用 \(\cos z=(e^{iz}+e^{-iz})/2\) 解出 \(e^{iz}=2\pm\sqrt3\)，取对数得

\[
z=\pm i\ln(2+\sqrt3)=\pm i\operatorname{arcosh}2\approx\pm 1.32\,i
\]

（再加上 \(2\pi\) 的整数平移）。这些是极点，到实轴的距离是 \(1.32\)。

\(g\) 的分母在 \(\sin z=\pm i/5\) 处为零，即 \(z=\pm i\operatorname{arsinh}(1/5)\approx\pm 0.199\,i\)。距离 \(0.199\)。这个数字并非巧合：Runge 函数 \(1/(1+25x^2)\) 的极点就在 \(\pm i/5\)，周期化只是把它折了个弯，分母里的 \(25\) 原封不动地留了下来。

\begin{center}
\begin{tikzpicture}[scale=0.85,font=\small,>=stealth]
  \begin{scope}
    \draw[fill=blue!10,draw=none] (-2.5,-1.58) rectangle (2.5,1.58);
    \draw[fill=orange!25,draw=none] (-2.5,-0.24) rectangle (2.5,0.24);
    \draw[->,black!55] (-2.7,0) -- (2.95,0) node[right,font=\footnotesize] {\(\operatorname{Re}z\)};
    \draw[->,black!55] (0,-2.0) -- (0,2.15) node[above,font=\footnotesize] {\(\operatorname{Im}z\)};
    \foreach \sx in {-1.5,1.5}{
      \foreach \sy in {-1.58,1.58}{
        \draw[blue!65!black,thick] (\sx-0.12,\sy-0.12) -- (\sx+0.12,\sy+0.12);
        \draw[blue!65!black,thick] (\sx-0.12,\sy+0.12) -- (\sx+0.12,\sy-0.12);
      }
      \foreach \sy in {-0.24,0.24}{
        \draw[orange!80!black,thick] (\sx-0.09,\sy-0.09) -- (\sx+0.09,\sy+0.09);
        \draw[orange!80!black,thick] (\sx-0.09,\sy+0.09) -- (\sx+0.09,\sy-0.09);
      }
    }
    \node[font=\footnotesize,blue!55!black] at (0.72,1.86) {\(a=1.32\)};
    \node[font=\footnotesize,orange!70!black] at (3.5,0.85) {\(a=0.20\)};
    \draw[orange!70!black] (3.0,0.72) -- (2.1,0.24);
    \node[font=\footnotesize,black!70] at (0,-2.35) {奇点到实轴的距离};
  \end{scope}
  \begin{scope}[shift={(7.6,1.7)}]
    \draw[->,black!55] (0,0) -- (5.7,0) node[right,font=\footnotesize] {\(k\)};
    \draw[black!55] (0,0.18) -- (0,-3.75);
    \draw[blue!65!black,thick] (0,0) -- (4.07,-3.5);
    \draw[orange!80!black,thick] (0,0) -- (5.0,-0.65);
    \node[right,font=\footnotesize,orange!70!black] at (5.05,-0.62) {斜率 \(-a\)};
    \node[right,font=\footnotesize,blue!55!black] at (3.35,-3.28) {斜率 \(-a\)};
    \foreach \yy/\lab in {0/{\(10^{0}\)},-1.75/{\(10^{-7}\)},-3.5/{\(10^{-14}\)}}{
      \draw[black!45] (-0.1,\yy) -- (0.1,\yy);
      \node[left,font=\footnotesize,black!70] at (-0.14,\yy) {\lab};
    }
    \foreach \xx/\lab in {1.67/{10},3.33/{20},5/{30}}{
      \draw[black!45] (\xx,0.1) -- (\xx,-0.1);
      \node[above,font=\footnotesize,black!70] at (\xx,0.12) {\lab};
    }
    \node[font=\footnotesize,black!70] at (2.8,-4.05) {\(\abs{\widehat{f}(k)}\)（纵轴对数刻度）};
  \end{scope}
\end{tikzpicture}

\emph{左边是两个函数的极点位置，右边是它们的 Fourier 系数。左图那条窄橙带把右图那条平缓的直线锁死了：极点离实轴 \(a\)，系数就以 \(e^{-a\abs{k}}\) 衰减，别的都不参与。}
\end{center}

右图两条直线的斜率，就是左图两个距离。\(1.32\) 与 \(0.199\) 之比约为 \(6.6\)，正好是开头那个``差出六七倍项数''的来源。要预测一个函数的谱精度，不必看它在实轴上有多光滑，要看它离开实轴多远才撞上第一个奇点。

下面把这句话变成推导。

\subsection{解析这个词比无穷可微严得多}

复可微在一点 \(z_0\) 存在，要求极限

\[
f'(z_0)=\lim_{h\to 0}\frac{f(z_0+h)-f(z_0)}{h}
\]

与 \(h\) 沿复平面哪个方向趋于零无关。实可微只管左右两个方向对得上，复可微要求无穷多个方向全部对得上。区域内处处复可微的函数称为解析函数。

这一步收紧带来的后果比预期猛烈。解析函数自动无穷次可微，而且它的 Taylor 级数在收敛圆盘内真的等于函数本身，不只是形式上凑出一个级数。实分析里那个标准反例 \(e^{-1/x^2}\)——所有导数在原点为零，函数却不恒为零——在复平面上无法存在：若它在原点解析，幂级数全零，那它在一个有聚点的集合（实轴上的一段区间）上取零；恒等定理随即迫使它在整个连通解析域上恒为零，与假设矛盾。

所以 \(e^{-1/x^2}\) 在实轴上无穷可微，却不能延拓成任何一个包含原点的复邻域上的解析函数。这类函数的 Fourier 系数衰减得比任何多项式都快，又慢于任何指数——它们卡在中间，因为它们的``条带宽度''是零。

恒等定理还有一个正面用途：若两个解析函数在一个有聚点的集合上相等，它们在共同的连通解析域上处处相等。因此实轴上给定的函数，解析延拓一旦存在就是唯一的。往复平面走没有自由度，只有一条路，走到奇点为止。

\subsection{条带有多宽，系数就衰减多快}

设 \(f\) 以 \(2\pi\) 为周期，且能解析延拓到条带 \(\abs{\operatorname{Im}z}<a\) 并在其中有界。Fourier 系数

\[
\widehat{f}(k)=\frac1{2\pi}\int_0^{2\pi}f(x)e^{-ikx}\,dx
\]

写在实轴上，但被积函数已经不只属于实轴了。

取 \(k>0\)，把积分路径从实轴整体下移到 \(\operatorname{Im}z=-b\)，其中 \(0<b<a\)。移动的合法性来自 Cauchy 积分定理：两条水平线段加上两端的竖直线段围成一个矩形，矩形内无奇点，环路积分为零；而周期性使两条竖直线段上的被积函数逐点相同、方向相反，抵消掉。剩下的等式说，两条水平线段上的积分相等。

新路径上 \(z=x-ib\)，于是

\[
\begin{aligned}
e^{-ikz}&=e^{-ik(x-ib)}=e^{-ikx}e^{-kb},\\
\abs{\widehat{f}(k)}&\le\frac1{2\pi}\int_0^{2\pi}\abs{f(x-ib)}\,e^{-kb}\,dx
\le\Bigl(\sup_{\abs{\operatorname{Im}z}\le b}\abs{f(z)}\Bigr)e^{-b\abs{k}}.
\end{aligned}
\]

\(k<0\) 时路径上移，得到同样的因子。令 \(b\) 趋近 \(a\)，就得到

\[
\abs{\widehat{f}(k)}\le C_b\,e^{-b\abs{k}},\qquad \forall\,b<a.
\]

条件对账值得单列。推导用到三件事：条带内无奇点（否则 Cauchy 定理失效，得改用留数，那是下一节的事）；\(f\) 在条带内有界（否则 \(C_b\) 发散，指数因子被吃掉）；周期性（否则竖直端段不抵消，需要另行估计它们在无穷远处的贡献）。三条缺一，结论就不成立或需要改写。而结论里的 \(b\) 只能取到 \(a\) 以下——常数 \(C_b\) 通常随 \(b\to a\) 一起爆掉，这是速率与常数之间的老式交换。

回头看开头那两个函数：\(f\) 的 \(a=\operatorname{arcosh}2\approx1.32\)，\(g\) 的 \(a=\operatorname{arsinh}(1/5)\approx0.199\)。图里那两条斜率不是拟合出来的。

\subsection{Runge 现象的病根写在 \(\pm i/5\) 上}

同一条原理换个场合，能解释一件看上去完全无关的事。

\hyperref[sec:08-Numerical-Analysis/03-Interpolation-and-Approximation/02]{``Runge 现象与节点选择''}里那个反例是 \(f(x)=1/(1+25x^2)\) 在 \([-1,1]\) 上的等距高次插值：节点越密，两端振荡越猛，最大误差随次数指数增长。当时的解释停在 Lebesgue 常数上——等距节点把最优逼近误差放大了指数倍。可放大因子为什么偏偏在这个函数上要了命，而在别的同样光滑的函数上不要命，那一节没能说透。

复平面给出答案。在原点做 Taylor 展开，\(1/(1+25x^2)=1-25x^2+625x^4-\cdots\) 的收敛半径是 \(1/5\)，尽管函数在整条实轴上有界且无穷可微。半径不是被实轴上的任何特征限制的，是被 \(\pm i/5\) 那对极点限制的：幂级数的收敛圆必须止步于最近的奇点，而最近的奇点根本不在实轴上。把 \(25\) 换成 \(1\)，极点退到 \(\pm i\)，收敛半径变成 \(1\)——实轴上的图像几乎没变，收敛行为却整个换了。

插值也一样。多项式插值的误差可以写成一个围道积分，被积核里带着节点多项式 \(\omega_n(x)/\omega_n(z)\) 的比值；插值收敛与否，取决于 \(f\) 的奇点落在 \(\abs{\omega_n}\) 的哪条等值线上。对 \([-1,1]\) 上的等距节点，这族等值线是一圈卵形，临界那条与虚轴交于约 \(\pm0.53i\)。\(1/(1+x^2)\) 的极点 \(\pm i\) 在这条卵形之外，等距插值收敛；\(1/(1+25x^2)\) 的极点 \(\pm 0.2i\) 落在里面，于是发散。两个函数在实轴上的差别只是一个胖瘦，判决却写在复平面上。

Chebyshev 节点把这族等值线换成了以 \([-1,1]\) 为焦点连线的椭圆。任何不在区间上的奇点，总能被某个椭圆包在外面，所以 Chebyshev 插值对区间邻域内解析的函数总是收敛，速率由穿过最近奇点的那个椭圆的参数 \(\rho\) 给出，误差按 \(\rho^{-n}\) 下降。这与本节的条带宽度不是两件事：Joukowski 映射 \(x=(w+w^{-1})/2\) 把椭圆族拉直成同心圆环，也就把 Chebyshev 系数的衰减翻译成了 Fourier 系数的衰减，\(\ln\rho\) 扮演的正是 \(a\) 的角色。

\hyperref[sec:08-Numerical-Analysis/08-Fourier-and-Spectral-Methods/04]{``谱方法：用光滑性换高精度''}里那句``解析函数常表现为几何收敛''，到这里可以换成一句能算的：几何收敛的公比等于 \(e^{-a}\)，\(a\) 是最近奇点到实轴（或到区间）的距离。要提高谱方法的精度，除了加项数，还有一条路是换变量把奇点推远——这正是各类坐标映射技巧的全部意图。

\subsection{条带宽度只给下界}

反过来读要小心。定理说的是``条带宽 \(a\) 保证衰减至少 \(e^{-b\abs{k}}\)''，没说``衰减恰好这么快''。整函数在全平面无奇点，条带可以任意宽，于是衰减超过任何指数；但具体多快，取决于函数离开实轴之后长得多猛。Paley--Wiener 理论把这件事说清楚了：整函数的增长阶与 Fourier 变换支集的形状精确对应，指数型 \(\sigma\) 的平方可积整函数对应支集在 \([-\sigma,\sigma]\) 内的变换。

采样定理站的也是这条线。严格带限信号按 Paley--Wiener 就是这样一个指数型整函数，全平面无奇点——可它的时域尾巴并不快：理想低通对应的 sinc 函数只按 \(1/\abs{t}\) 代数衰减。带宽管住的是采样率的下限（频谱副本不重叠），管不了时域衰减；两者之间隔着的正是奇点分布这一层。\hyperref[sec:08-Numerical-Analysis/08-Fourier-and-Spectral-Methods/02]{``采样、混叠与可恢复性''}里那句``理想矩形滤波器的时域响应正是无限支撑 sinc''，说的就是这件事的工程后果。

本节全程只做了一个动作：把积分路径从实轴挪开，而条带内没有奇点，所以什么代价都不用付。下一节把路径继续挪，一直挪到它绕过了奇点——那时代价出现了，并且可以精确计算。
